\documentclass[12pt, a4paper]{exam}
\usepackage{fontspec}
\usepackage{xunicode}
\usepackage{xltxtra}
\usepackage{polyglossia}
\usepackage{graphicx}
\usepackage{tikz}
\usetikzlibrary{positioning, calc}
\usepackage{geometry}
\usepackage{enumitem}
\usepackage{multicol}
\usepackage{array}
\usepackage{lastpage}
\usepackage{amssymb}
\usepackage{multirow}
\usepackage{amsmath}

\setmainlanguage{thai}
\XeTeXlinebreaklocale "th"
\XeTeXlinebreakskip = 0pt plus 1pt

% Page Margins
\geometry{top=0.8in, bottom=1in, left=0.8in, right=0.8in}

% Fonts
\newfontfamily\thaifont[
  Scale=1.33,
  Script=Thai,
  Path = ../fonts/,
  Extension = .ttf,
  BoldFont = *-Bold,
  ItalicFont = *-Italic,
  BoldItalicFont = *-BoldItalic
]{THSarabunNew}

\setmainfont[
  Scale=1.33,
  Path = ../fonts/,
  Extension = .ttf,
  BoldFont = *-Bold,
  ItalicFont = *-Italic,
  BoldItalicFont = *-BoldItalic
]{THSarabunNew}

% Header and Footer
\pagestyle{headandfoot}
\extraheadheight{0.2in}

\newcommand{\examheader}{%
    ชื่อ-นามสกุล \makebox[3.0in]{\hrulefill} รหัสนักศึกษา \makebox[2.0in]{\hrulefill}
}

\firstpageheader{\examheader}{}{}
\runningheader{\examheader}{}{}

\firstpagefooter{\textit{MTH234 Question Bank (Full Version)}}{}{\thepage}
\runningfooter{\textit{MTH234 Question Bank (Full Version)}}{}{\thepage}

% Custom Thai Choice Labels
\makeatletter
\newcommand{\thaialph}[1]{%
  \ifcase\value{#1}\or ก\or ข\or ค\or ง\or จ\else\@ctrerr\fi
}
\makeatother

% Question Format
\renewcommand{\thequestion}{\arabic{question}}
\renewcommand{\questionlabel}{\thequestion.}

% Choice Format
\renewcommand{\thechoice}{\thaialph{choice}}
\renewcommand{\choicelabel}{\thechoice)}

\renewcommand{\choiceshook}{%
  \setlength{\leftmargin}{2.5em}%
  \setlength{\labelwidth}{20pt}%
  \setlength{\labelsep}{5pt}%
}

\begin{document}

\section*{Part 1: Assignment 1 (4 Questions)}
\begin{questions}

% Assignment Q1
\question กำหนดให้ $A = \begin{bmatrix} 1 & -2 \\ 4 & 3 \\ -1 & 1 \end{bmatrix}$
\begin{parts}
    \part จงหาเมทริกซ์ $X$ ที่ทำให้ $2A - 3X = 4(A + X)$
    \part จงหาค่า $AA^T$ (โดยที่ $A^T$ คือทรานสโพสของ $A$)
\end{parts}
\vspace{3cm}

% Assignment Q2
\question กำหนดให้
\[
A = \begin{bmatrix} 1 & -3 & 2 & 1 \\ 3 & -9 & 10 & 2 \\ 2 & -6 & 4 & 2 \\ 2 & -6 & 8 & 1 \end{bmatrix}
\]
จงใช้ลำดับของการดำเนินการตามแถวเบื้องต้น (elementary row operations) เพื่อหาเมทริกซ์ที่เป็นรูปแบบขั้นบันไดตามแถวแบบลดรูป (reduced row echelon form) ซึ่งสมมูลกันทางแถวกับเมทริกซ์ $A$
\vspace{5cm}

% Assignment Q3
\question กำหนดให้ $A = \begin{bmatrix} 2 & 6 & 1 \\ 3 & 9 & 2 \\ 0 & -1 & 3 \end{bmatrix}$ จงใช้การดำเนินการตามแถวเบื้องต้นเพื่อหาค่าดีเทอร์มิแนนต์ของ $A$ ($\det(A)$) และเมทริกซ์ผกผัน $A^{-1}$
\vspace{5cm}

% Assignment Q4
\question พิจารณาระบบสมการเชิงเส้น:
\[
\begin{aligned}
x + 2y + 5z &= 15 \\
-x + 7y - 5z &= -20 \\
-x + 8y + 3z &= 4
\end{aligned}
\]
จงหาผลเฉลยโดย:
\begin{parts}
    \part ใช้การดำเนินการตามแถวเบื้องต้น
    \part ใช้กฎของคราเมอร์
\end{parts}
\vspace{5cm}

\end{questions}

\filbreak
\section*{Part 2: Mock Exam 1 (6 Questions)}
\begin{questions}

% Mock1 Q1
\question (6 คะแนน) จงแก้สมการเมทริกซ์ต่อไปนี้เพื่อหาค่า $X$:
\[
X^T \begin{bmatrix} 3 & 1 \\ -1 & 2 \end{bmatrix} - \begin{bmatrix} 1 & 4 \\ 2 & 0 \end{bmatrix} = \begin{bmatrix} 2 & -2 \\ 5 & 4 \end{bmatrix}
\]
โดยที่ $X^T$ คือทรานสโพสของเมทริกซ์ $X$
\vspace{5cm}

% Mock1 Q2
\question (10 คะแนน) จงใช้การดำเนินการตามแถวเบื้องต้นเพื่อหา $\det(A)$ และ $A^{-1}$ (ถ้ามี) เมื่อ
\[
A = \begin{bmatrix} 0 & 0 & 2 & 0 \\ 1 & 0 & 0 & 1 \\ 0 & -1 & 3 & 0 \\ 2 & 1 & -5 & 3 \end{bmatrix}
\]
\vspace{5cm}

% Mock1 Q3
\question (8 คะแนน) จงใช้กฎของคราเมอร์ (Cramer's Rule) เพื่อแก้ระบบสมการต่อไปนี้:
\[
\begin{aligned}
2x_1 + x_2 - x_3 &= 4 \\
x_1 - 2x_2 + x_3 &= 1 \\
3x_1 - 3x_2 + 2x_3 &= 0
\end{aligned}
\]
\vspace{5cm}

% Mock1 Q4
\question (8 คะแนน) จงใช้การดำเนินการตามแถวเบื้องต้นเพื่อตรวจสอบว่าระบบสมการเชิงเส้นต่อไปนี้มีคำตอบ (consistent) หรือไม่ หากมี จงหาผลเฉลยทั่วไป:
\[
\begin{aligned}
x_1 - 3x_2 + x_3 &= 2 \\
3x_1 - 8x_2 + 2x_3 &= 5 \\
2x_2 - 2x_3 &= -2
\end{aligned}
\]
\vspace{5cm}

% Mock1 Q5
\question (8 คะแนน) จงหาค่าของ $a$ ทั้งหมดที่ทำให้ระบบสมการเชิงเส้น:
\[
\begin{aligned}
x + y - z &= 1 \\
2x + 3y + az &= 3 \\
x + ay + 3z &= 2
\end{aligned}
\]
อยู่ในแต่ละกรณีต่อไปนี้:
(a) ไม่มีผลเฉลย
(b) มีผลเฉลยชุดเดียว
(c) มีผลเฉลยมากมายไม่จำกัด
\vspace{5cm}

% Mock1 Q6
\question (10 คะแนน) สำหรับระบบสมการเชิงเส้น:
\[
\begin{aligned}
w + 2x + 2y + 2z &= 1 \\
w + 3x + 4y + 3z &= 6 \\
2w + 6x + 9y + 8z &= 13 \\
w + 5x + 8y + \omega z &= \varepsilon
\end{aligned}
\]
จงหาค่าของ $\omega$ และ $\varepsilon$ ที่ทำให้ระบบสมการมีคำตอบในกรณีต่อไปนี้:
(a) กรณี $\omega = 6, \varepsilon = 17$
(b) กรณี $\omega = 5, \varepsilon = 16$
\vspace{5cm}

\end{questions}

\filbreak
\section*{Part 3: Mock Exam 2 (7 Questions)}
\begin{questions}

% Mock2 Q1
\question (10 คะแนน) กำหนดให้ $A = \begin{bmatrix} -2 & 0 & 4 \\ 2 & 6 & 3 \\ 4 & -1 & 2 \end{bmatrix}$, $B = \begin{bmatrix} 1 & 0 & -3 \\ -2 & -5 & 0 \\ -2 & 1 & 2 \end{bmatrix}$ และ $C = \begin{bmatrix} 0 & -2 & 0 \\ -1 & -3 & 4 \end{bmatrix}$ จงหา:
\begin{parts}
    \part $3A - B^T$
    \part $BC$ และ $CB$
    \part เมทริกซ์ $X$ ที่ทำให้ $X(A+B) = C$
\end{parts}
\vspace{5cm}

% Mock2 Q2 (Duplicate type as Assign Q2 but different matrix)
\question (12 คะแนน) กำหนดให้
\[
A = \begin{bmatrix} 1 & 0 & -2 & 1 & 0 \\ 0 & -1 & -3 & 1 & 3 \\ -2 & -1 & 1 & -1 & 3 \\ 0 & 3 & 9 & 0 & -12 \end{bmatrix}
\]
จงหา Reduced Row Echelon Form (RREF) ของเมทริกซ์ $A$
\vspace{5cm}

% Mock2 Q3
\question (15 คะแนน) จงหา $\det(A)$ และ $A^{-1}$ โดยใช้ Row Operations:
\[
A = \begin{bmatrix} 2 & 6 & 6 & 2 \\ 2 & 7 & 3 & 6 \\ 1 & 5 & 0 & 1 \\ 3 & 7 & 0 & 7 \end{bmatrix}
\]
\vspace{5cm}

% Mock2 Q4 (Same as Assign Q4 but diff numbers)
\question (13 คะแนน) พิจารณาระบบสมการเชิงเส้น:
\[
\begin{aligned}
2x_1 + 3x_2 + 5x_3 &= 4 \\
3x_1 + 5x_2 + 9x_3 &= 7 \\
5x_1 + 9x_2 + 13x_3 &= 17
\end{aligned}
\]
จงหาคำตอบโดยใช้ (a) Elementary Row Operations และ (b) Cramer's Rule
\vspace{5cm}

% Mock2 Q5
\question (13 คะแนน) จงตรวจสอบว่า $A\mathbf{x} = \mathbf{b}$ มีผลเฉลยที่ไม่ชัด (non-trivial solutions) หรือไม่ และหาผลเฉลยทั่วไปและฐาน (basis) ของปริภูมิผลเฉลย:
\[
A = \begin{bmatrix} 1 & 2 & 3 \\ 2 & -1 & 1 \\ 3 & 1 & 4 \\ 5 & 1 & 6 \end{bmatrix}, \quad \mathbf{b} = \begin{bmatrix} 0 \\ 0 \\ 0 \\ 0 \end{bmatrix}
\]
\vspace{5cm}

% Mock2 Q6 (Similar to Mock1 Q5)
\question (12 คะแนน) จงหาค่า $k$ ที่ทำให้ระบบสมการ:
\[
\begin{aligned}
x + y + kz &= 3 \\
x + ky + z &= 2 \\
kx + y + z &= 1
\end{aligned}
\]
มี (a) ไม่มีผลเฉลย, (b) มีผลเฉลยเดียว, (c) มีหลายผลเฉลย
\vspace{5cm}

% Mock2 Q7 (Extra)
\question (15 คะแนน) จงหาค่าดีเทอร์มิแนนต์ของเมทริกซ์ต่อไปนี้โดยใช้สมบัติของดีเทอร์มิแนนต์:
\[
A = \begin{bmatrix} 1 & -2 & -1 & -2 & 1 \\ 2 & 1 & -2 & 1 & -3 \\ 3 & -1 & -3 & -1 & -1 \\ 4 & 1 & 1 & 0 & 1 \\ 5 & -5 & -5 & 2 & 0 \end{bmatrix}
\]
\vspace{5cm}

\end{questions}
\end{document}
